\addcontentsline{toc}{section}{Введение}
\section*{Введение}

В ходе выполнения лабораторной работы нужно реализовать двумерный клеточный автомат с окрестностью фон Неймана в соответствии с полученным номером №. Граничный условия - торроидальные. Пользователь определяет ширину поля и количество итераций. Необходимо учесть возможность ввода различных начальных условий (как вручную, так и случайным образом) по выбору пользователя. Реализация возможна в консоли или в графике. В итоге необходимо проанализировать свой клеточный автомат в отчете (его поведение, паттерны, "сходимость" и т.д.). 

Функция изменения состояния клеток расчитывается по формуле:

\[
№ = \text{номер варианта} * 11 * \text{год рождения} * \text{день} * \text{месяц}
\]